% The screening framework. Brief: briefs/framework.md. Claims carried: the entries of ../claims.yaml with section
% framework (CL-M18, CL-M19), with the design and methods entries it summarises. Prose: PH-W8.13. It states the inputs,
% the stage logic, the outputs and how to add a candidate phase, in about 600 words and one compact table. No number is
% typed and no result is stated; the release location is a pending field the authors fill (\nacreRepoURL).
\section{The screening framework}\label{sec:framework}

% (1) What the framework does and its inputs.
The framework turns sourced thermochemistry into stage-resolved verdicts for any candidate mineral phase (Table~\ref{tab:framework}). % CL-M18
It has three inputs. % CL-M18
The thermochemical data are solubility products, formation free energies, enthalpies and entropies, molar volumes and interfacial energies, each with its uncertainty. % CL-M18, CL-M11
Every value carries its DOI, its locator and an evidence grade, and the grade limits the prior built from it (Section~\ref{sec:methods}). % CL-M11, CL-M01
The candidate list names each phase with its formula unit, its stage and its framing, crystalline or amorphous. % CL-M18
The windows are the grid of pH, calcium dose and anion dose at the processing stage and the cutoff classes of the cell stage. % CL-M02
An input with no sourced value leaves its test undetermined, and an undetermined test never passes. % CL-M11, CL-D09

\begin{table}[t]
\small
\caption{The screening framework: inputs, stage logic, outputs and how a phase is added. Symbols as in Section~\ref{sec:design}; file paths are those of the released repository.}
\label{tab:framework}
\begin{tabular*}{\columnwidth}{@{}>{\raggedright\arraybackslash}p{0.21\columnwidth}>{\raggedright\arraybackslash}p{0.75\columnwidth}@{}}
\hline
Part & Content\\
\hline
Inputs & sourced thermochemistry with uncertainties and evidence grades; the candidate list; the pH and dose grid and the cutoff classes\\
Processing stage & dissolution test, $\Omb\ge1$ at every state of the grid ($K_2$); nucleation and thin-film tests ($K_3$, $K_4$)\\
Cell stage & upper bound $E_\mathrm{ub}$ of the oxidation onset for a saturated sink; exclusion per cutoff class ($K_1$)\\
Between stages & water-activity thresholds of the conversion steps\\
Outputs & exclusion map; break-even ratio $r^*$; required-precision map; the input value at which each verdict flips\\
New phase & data in \texttt{lit/params.yaml}; phase in \texttt{configs/preprint/screen.yaml}; reactions in the onset ladder; rerun the screen\\
\hline
\end{tabular*}
\end{table}

% (2) The stage logic.
Each phase is tested at the stage at which it would exist. % CL-D11
At the processing stage the screen asks whether the phase dissolves in the slurry. % CL-D05
It computes the supersaturation $\Omb$ of Eq.~\eqref{eq:stability} at every state of the grid. % CL-D05
A phase passes only over a region in which every state has $\Omb\ge1$, and that pass is a necessary condition, not a sufficient one. % CL-D05, CL-D20
At the cell stage the screen bounds the oxidation onset from above. % CL-D08
Releasing the cation into a solid salt of the electrolyte anion gives the upper bound $E_\mathrm{ub}$ for a saturated electrolyte. % CL-D08
A phase whose upper bound plus its margin lies below a cutoff is thermodynamically excluded\nacreExcludedMark\ at that cutoff. % CL-D08
A pass would need a lower bound, and the lower bound needs transfer free energies that are unmeasured. % CL-M17
The cell stage can therefore exclude a phase but cannot admit one. % CL-M17
Conversion steps between the stages are water-activity thresholds only (Section~\ref{sec:design}). % CL-D11
Kinetic stabilisation lies outside the screen; only a measured service limit, tier-two evidence, can show it. % CL-D09

% (3) The outputs.
The framework returns four outputs. % CL-M18
The exclusion map gives each phase a verdict per stage and per cutoff class (Table~\ref{tab:regime}). % CL-S13
The break-even ratio $r^*$ states how far below saturation the electrolyte may fall before an exclusion lapses (Table~\ref{tab:breakeven}). % CL-S03
The required-precision map gives the precision a measurement needs to decide each open verdict (Section~\ref{sec:res-precision}). % CL-P04
The flip values give the input value at which each verdict would change. % CL-P04
The relations are written per formula unit for any mineral, so a new phase receives the same four outputs. % CL-R09, CL-M18

% (4) How to add a candidate phase.
A new phase needs a sourced solubility product with its uncertainty for the dissolution test. % CL-M18
It needs the formation thermochemistry of the phase and of its oxidation products for the upper bound. % CL-M18
The nucleation and thin-film tests also need its molar volume and interfacial energies; without them those tests read undetermined. % CL-M18, CL-D09
The values go into the parameter file, \texttt{lit/params.yaml}, each with its DOI, locator and grade. % CL-M18
The phase goes into the candidate list, \texttt{configs/preprint/screen.yaml}, with its formula unit, stoichiometry and parameter keys. % CL-M18
Its oxidation reactions go into the onset ladder of the package \texttt{nacre.pourbaix}. % CL-M18
The screen is then rerun through \texttt{python -m nacre.screen}, and \texttt{nacre preprint regime-table} rebuilds the map from the outputs. % CL-M18
Every number of the rebuilt map is traced to its source file. % CL-M09

% (5) Release.
The code and data are released as the open repository \texttt{biomin-screen} (\nacreRepoURL). % CL-M19
